\documentclass[12pt]{article}
\usepackage[utf8]{inputenc}
\usepackage[T1]{fontenc}
\usepackage{hyperref}
\title{Dakka Language Resources for K-12 Curriculum Development}
\author{Compiled by Claude Sonnet 5.5 (claude-sonnet-5-5) for Andr\'as Kornai}
\date{}
\begin{document}
\maketitle

\section{Introduction}
Dakka (ISO 639-3 \texttt{dkk}; Glottocode dakk1238) is an endangered Austronesian language of Sulawesi, Indonesia, spoken in the Wonomulyo district of Polewali Mandar Regency. It belongs to the Pitu Ulunna Salu subgroup of Northern South Sulawesi, together with Bambam, Aralle-Tabulahan, Ulumanda' and Pannei. Wikipedia lists about 1,000 speakers (2010). Almost nothing educational was found.

\section{Primary Educational Resources}
No primers, textbooks or school programmes for Dakka were found in search results.

\section{Linguistic and Reference Materials}
\begin{itemize}
\item Wikipedia, Dakka language: \url{https://en.wikipedia.org/wiki/Dakka_language}
\item Glottolog record: \url{https://glottolog.org/resource/languoid/id/dakk1238.bigmap.html}
\item Wikipedia, Bambam language (a closely related language, not Dakka): \url{https://en.wikipedia.org/wiki/Bambam_language}
\item The languages of Sulawesi (ANU Open Research repository, general survey): \url{https://openresearch-repository.anu.edu.au/items/0d73a31f-b139-4c9a-8063-8e2bc9e6f10f}
\end{itemize}

\section{Religious and Community Resources}
Joshua Project has a profile of the Dakka people in Indonesia (reported population about 22,000); its cultural claims appear generic and should be treated with caution: \url{https://joshuaproject.net/people_groups/11476}

\section{Digital and Audio Resources}
None found.

\section{Curriculum Development Context}
No evidence of Dakka teaching in schools was found. Since the language is spoken mainly by adults, any curriculum would be a revitalisation effort needing fieldwork and documentation first.

\end{document}